\section{Summary and Future Directions}
\label{sec:pho_summary}

This chapter presented an object-centric force-policy framework for pushing heavy rigid objects along prescribed paths under unknown mass and friction. The policy is trained in simulation using force-space actions and a descriptor that summarizes contact resistance. At deployment, the robot estimates this descriptor through a short probing interaction and executes the requested forces through compliant control. Experiments on a Boston Dynamics Spot mobile manipulator demonstrated deployment of the simulation-trained policy without policy retraining.

The framework extends the interaction-abstraction perspective from \cref{ch:flex}. Whereas \FLEX{} exploits the geometric structure of articulated objects, this chapter represents rigid-object motion relative to a desired path and conditions the policy on contact resistance. In both cases, force space separates the object-level interaction from robot-specific control. In the terminology of \cref{ch:novelty_open_world}, the two frameworks expand the range of variation that can be assimilated within an existing policy representation: the evaluated changes in objects, physical conditions, and embodiments do not require the policy itself to be revised.

The interaction descriptor should be interpreted as a coarse regime indicator rather than a precise estimate of mass or friction: replacing it with its in-distribution mean produced performance close to that of the full system, while extreme or absent descriptor values degraded tracking. The generalization demonstrated here therefore arises from the combination of object-centric force-space representation, training across physical variation, and coarse online conditioning, rather than from exact physical-parameter identification.

Several limitations remain. The framework assumes quasi-static planar transport and conditions on a single scalar contact-resistance descriptor. This is sufficient for the evaluated setting but cannot represent spatially varying friction, anisotropic contact, time-varying resistance, or dynamic manipulation regimes. Future work should extend the descriptor to richer learned representations, incorporate contact-state recovery when contact is lost, and support multi-contact or collaborative manipulation where multiple forces must be coordinated across contact points.

Together, \cref{ch:flex,ch:pho} provide the generalization counterpart to the accommodation studied in \cref{ch:rapid_learn}. Chapter~\ref{ch:rapid_learn} showed how action abstractions localize the behavior that must be learned after an execution failure. Chapters~\ref{ch:flex} and~\ref{ch:pho} showed how interaction abstractions allow a fixed policy to remain applicable across the evaluated variation in objects, physical conditions, and embodiments. The contributions therefore support the two sides of the dissertation's technical argument: structured abstractions can localize revision when existing competence is insufficient and can reduce the need for revision by capturing invariances that support generalization.

\section{Technical Appendix}
\label{sec:pho_technical_appendix}

\subsection{Physics Engine and Scene Configuration}

All simulation experiments are conducted in MuJoCo~\cite{todorov2012mujoco}.
The scene consists of a ground plane (half-size $5 \times 5 \times 0.1$\,m) and a single rigid box body with a free joint (6-DoF).
Table~\ref{tab:mujoco_physics} lists the relevant physics and contact parameters.

\begin{table}[h]
\centering
\caption{MuJoCo physics and contact parameters.}
\label{tab:mujoco_physics}
\begin{tabular}{ll}
\toprule
\textbf{Parameter}              & \textbf{Value}              \\
\midrule
Integration timestep            & $0.0025$\,s                 \\
Substeps per RL step            & $10$                        \\
Effective control $\Delta t$    & $0.025$\,s                  \\
Contact dimension (\texttt{condim}) & $4$ (tangent + torsional) \\
\texttt{solimp}                 & $[0.9, 0.95, 0.001, 0.5, 2]$ \\
Linear damping (box DoFs)       & $0.05$                      \\
Angular damping (box DoFs)      & $0.05$                      \\
Warm-start settle steps         & $50$ ($0.125$\,s)           \\
\bottomrule
\end{tabular}
\end{table}

\subsection{Object Geometry}

The box geometry is specified as half-extents in the MuJoCo model.
The full dimensions are $0.465 \times 0.61 \times 0.63$\,m (width $\times$ depth $\times$ height), matching the real-world cardboard box used in hardware experiments.
The nominal inertial properties (overridden at each episode by domain randomization) are: mass $15.0$\,kg, diagonal inertia $[0.2403, 0.1916, 0.1839]$\,kg\,m$^2$.
The box is initialized at height $z = 0.315$\,m (half the box height) at the first waypoint of the reference path.

\subsection{Reference Path Generation}

During training, the reference trajectory is a circular arc with radius $1.5$\,m, spanning $120^\circ$ (from $-\pi/3$ to $\pi/3$), discretized into $50$ equispaced waypoints.
Each training episode samples a short segment of arc length $0.3$\,m from this full arc.
The segment start position is sampled uniformly along the arc.

\subsection{Domain Randomization}

At each episode reset, the following physical parameters are independently sampled from uniform distributions:

\begin{itemize}
    \item \textbf{Object mass:} $m \sim \mathcal{U}(5, 15)$\,kg
    \item \textbf{Sliding friction coefficient:} $\mu \sim \mathcal{U}(0.4, 0.6)$
\end{itemize}

\noindent
The Coulomb friction coefficient is applied symmetrically to both the box and the floor geom.
Additional friction parameters (spinning friction $= 0.5$, rolling friction $= 0.01$) remain fixed.
Box dimensions are held constant during training; variation is introduced only at evaluation time.

\subsection{MDP Specification}
\label{app:mdp}

\subsubsection{State Space}

The state is a $6$-dimensional vector defined in the object-centric, path-relative frame:
%
\begin{equation}
    s_t = \bigl[\, e^{\mathrm{lat}}_t,\; e^{\theta}_t,\; v^{\parallel}_t,\; v^{\perp}_t,\; \omega_t,\; z \,\bigr] \in \mathbb{R}^6,
\end{equation}
%
where:
\begin{itemize}
    \item $e^{\mathrm{lat}}_t$: signed lateral deviation from the nearest path waypoint (m),
    \item $e^{\theta}_t$: wrapped orientation error relative to the local path tangent (rad),
    \item $v^{\parallel}_t$: velocity projected onto the path tangent (m/s),
    \item $v^{\perp}_t$: velocity projected onto the path normal (m/s),
    \item $\omega_t$: angular velocity about the vertical axis (rad/s),
    \item $z = \mathrm{clip}(\mu m g / F_{\max},\, 0,\, 1)$: normalized contact-resistance descriptor.
\end{itemize}

Velocities are computed via finite differences over $\Delta t = 0.025$\,s.
The interaction descriptor $z$ is computed from ground-truth physics in simulation and held constant within each episode.
In the push-from-edge configuration, the ``heading'' used for orientation error is the direction from the box centroid to the contact edge, expressed in the box body frame as $[0.465, 0.63, 0]$ (normalized).

\subsubsection{Action Space}

In the push-from-edge mode used for all reported experiments, the action space is $2$-dimensional.
The actor network outputs a unit direction vector $\hat{\mathbf{d}} \in \mathbb{R}^2$ and a scalar magnitude $m \in [0,1]$ (see Section~\ref{app:network}), combined as $\mathbf{a}_t = \hat{\mathbf{d}} \cdot m$.
The environment scales this to physical forces:
%
\begin{equation}
    F_x = \mathrm{clip}(a_x,\, {-1},\, 1) \cdot F_{\max}, \quad F_y = \mathrm{clip}(a_y,\, {-1},\, 1) \cdot F_{\max},
\end{equation}
%
where $F_{\max} = 200$\,N.
Torque arises implicitly from the lever-arm geometry: $\tau = \mathbf{r} \times \mathbf{F}$, where $\mathbf{r}$ is the vector from the box centroid to the contact edge point.
The direction-magnitude factorization constrains the action to lie within the unit disk, so the effective force magnitude is bounded by $F_{\max}$.

Forces are applied in the box body frame and rotated into the world frame at each physics substep.

\subsubsection{Reward Function}

The reward at each step is a sum of five terms:
%
\begin{equation}
    R = r_{\text{prog}} + r_{\text{speed}} + r_{\text{lat}} + r_{\text{ori}} + r_{\text{spin}},
\end{equation}
%
defined as follows.
Table~\ref{tab:reward} lists each component, its expression, and the corresponding weight.

\begin{table}[h]
\centering
\caption{Reward components. $\Delta t = 0.025$\,s; $v_{\mathrm{ref}} = 0.3$\,m/s.}
\label{tab:reward}
\small
\begin{tabular}{lll}
\toprule
\textbf{Term} & \textbf{Expression} & \textbf{Weight} \\
\midrule
Progress  & $v^{\parallel}_t \Delta t \cdot e^{-5|e^{\mathrm{lat}}_t|}$ & $+100$ \\
Speed     & $(v^{\parallel}_t - v_{\mathrm{ref}})^2$       & $-8$   \\
Lateral   & $(e^{\mathrm{lat}}_t)^2$                        & $-10$  \\
Orientation & $(e^{\theta}_t)^2$                             & $-2$   \\
Spin      & $\omega_t^2$                                    & $-0.5$ \\
\bottomrule
\end{tabular}
\end{table}

\noindent\textbf{Terminal rewards and penalties:}
\begin{itemize}
    \item \textbf{Success bonus:} $+100$ when progress $> 0.95$ and deviation $<$ goal threshold.
    \item \textbf{Failure penalty:} $-500$ when $|e^{\mathrm{lat}}_t| >$ deviation tolerance ($0.5$\,m).
\end{itemize}

The goal threshold is computed dynamically as $\max(0.10 \times L_{\mathrm{seg}},\, 0.03)$\,m, where $L_{\mathrm{seg}}$ is the current segment arc length.
Episodes are truncated after $500$ steps (short-segment training) or $1000$ steps (full-arc evaluation).

\subsection{TD3 Hyperparameters}
\label{app:td3}

We use Twin Delayed Deep Deterministic Policy Gradient (TD3)~\cite{fujimoto2018td3}.
Table~\ref{tab:td3_hyper} reports all hyperparameters.

\begin{table}[h]
\centering
\caption{TD3 training hyperparameters.}
\label{tab:td3_hyper}
\begin{tabular}{ll}
\toprule
\textbf{Hyperparameter}            & \textbf{Value}     \\
\midrule
Learning rate (actor \& critics)   & $1 \times 10^{-4}$ \\
Optimizer                          & Adam               \\
Discount factor ($\gamma$)         & $0.99$             \\
Target smoothing coefficient ($\tau$) & $0.995$         \\
Target policy noise ($\sigma$)     & $0.2$              \\
Target noise clip                  & $0.5$              \\
Policy delay                       & $2$                \\
Gradient steps per env step        & $2$                \\
Exploration noise ($\sigma_{\mathrm{expl}}$) & $0.1$    \\
Batch size                         & $256$              \\
Replay buffer capacity             & $500{,}000$        \\
Random exploration steps           & $2{,}000$          \\
Training episodes                  & $3{,}001$          \\
Evaluation frequency               & every $25$ episodes \\
Evaluation episodes per condition  & $25$               \\
Checkpoint save frequency          & every $500$ episodes \\
Seeds                              & $0$--$9$ ($10$ seeds total) \\
\bottomrule
\end{tabular}
\end{table}

\noindent\textbf{Replay buffer management.}
When the buffer exceeds its maximum capacity, the oldest $1/5$ of transitions are discarded.

\noindent\textbf{Model selection.}
The best model is selected every $25$ episodes on in-distribution conditions using the composite score
\[
0.5 \times \text{full-arc success rate}
+ 0.5 \times \text{generalization success rate}.
\]
The generalization evaluation uses random arc geometries with radius $r \sim \mathcal{U}(1, 2)$\,m and angular span variation of $\pm 60^\circ$ from the nominal $120^\circ$ training arc.

\subsection{Network Architecture}
\label{app:network}

\subsubsection{Actor}

The actor network takes the 6D state as input and produces a 2D force action.
The architecture uses a direction-magnitude decomposition:

\begin{enumerate}
    \item \textbf{Shared trunk:}\\
    $\mathrm{Linear}(6 \!\to\! 400) \!\to\! \mathrm{ReLU} \!\to\! \mathrm{Linear}(400 \!\to\! 300) \!\to\! \mathrm{ReLU}$
    \item \textbf{Direction head:}\\
    $\mathrm{Linear}(300 \!\to\! 2) \to \tanh \to \ell_2\text{-normalize}$
    \item \textbf{Magnitude head:}\\
    $\mathrm{Linear}(300 \!\to\! 1) \to \mathrm{sigmoid}$
\end{enumerate}

The output force is $\mathbf{f} = \hat{\mathbf{d}} \cdot m$, where $\hat{\mathbf{d}} \in \mathbb{R}^2$ is the unit direction vector and $m \in [0, 1]$ is the magnitude.
The direction-magnitude factorization ensures the force direction is always a valid unit vector, improving training stability.

\subsubsection{Critic}

Each of the two critic networks (twin critics) has the architecture:
%
\begin{align*}
    &\mathrm{Linear}(8 \to 400) \to \mathrm{ReLU} \\
    &\quad\to \mathrm{Linear}(400 \to 300) \to \mathrm{ReLU} \\
    &\quad\to \mathrm{Linear}(300 \to 1),
\end{align*}
%
where the input dimension is $8 = 6\text{ (state)} + 2\text{ (action)}$.

\noindent\textbf{Target policy smoothing.}
During target Q-value computation, Gaussian noise is added to the target actor's output.
The force component of the noisy action is re-normalized to a unit vector after noise addition to preserve the direction-magnitude structure.

\subsection{Simulation Evaluation Protocol}
\label{app:eval}

\subsubsection{In-Distribution Evaluation}

For each trained seed, we run $100$ evaluation rollouts per condition with a maximum of $1000$ steps per episode.
Physical parameters not under test are sampled from the in-distribution ranges ($m \sim \mathcal{U}(5, 15)$\,kg, $\mu \sim \mathcal{U}(0.4, 0.6)$) via the environment's domain randomization at each episode reset.
The evaluation seed is fixed at $42$ for reproducibility.
Results are reported as mean $\pm$ standard deviation of per-episode RMSE values across seeds.

Table~\ref{tab:ood_ranges} summarizes the out-of-distribution evaluation ranges.

\begin{table}[h]
\centering
\caption{Out-of-distribution evaluation ranges. When testing one parameter OOD, all others remain in-distribution.}
\label{tab:ood_ranges}
\begin{tabular}{lccc}
\toprule
\textbf{Parameter} & \textbf{ID Range} & \textbf{OOD Low} & \textbf{OOD High} \\
\midrule
Mass (kg)           & $[5, 15]$       & $[2.5, 5]$     & $[15, 17.5]$    \\
Friction ($\mu$)    & $[0.4, 0.6]$    & $[0.3, 0.4]$   & $[0.6, 0.7]$    \\
Path radius (m)     & $1.5$           & $[1.0, 1.5]$   & $[1.5, 2.0]$    \\
Box dims (cm)       & $46.5 \times 61 \times 63$ & $40^3$ & $70^3$     \\
\bottomrule
\end{tabular}
\end{table}

\subsubsection{Path Generalization Evaluation}

Path generalization is evaluated with all physical parameters in-distribution.
The tested trajectories are:
\begin{itemize}
    \item \textbf{Arc $60^\circ$, $120^\circ$, $180^\circ$:} Circular arcs of radius $1.5$\,m with the indicated angular span, centered at $0^\circ$.
    \item \textbf{S-path:} A sinusoidal path of total length $3.0$\,m and amplitude $0.5$\,m, discretized into $100$ waypoints.
    \item \textbf{Sine-wave (meandering):} A sinusoidal path of total length $4.0$\,m and amplitude $0.4$\,m, discretized into $150$ waypoints.
\end{itemize}

\subsubsection{RMSE Computation}

Tracking performance is measured using centroid-to-path RMSE:
%
\begin{equation}
    \mathrm{RMSE} = \sqrt{\frac{1}{N} \sum_{t=1}^{N} d_t^2},
\end{equation}
%
where $d_t = \min_{j} \| \mathbf{c}_t - \mathbf{q}_j \|_2$ is the Euclidean distance from the object centroid to the closest waypoint on the reference path, and $N$ is the total number of steps across all episodes in a given condition.

\subsubsection{Ablation: Interaction Descriptor}

The ablation study evaluates four descriptor settings under fully in-distribution conditions ($m \sim \mathcal{U}(5, 15)$\,kg, $\mu \sim \mathcal{U}(0.4, 0.6)$):
\begin{itemize}
    \item \textbf{Estimated (ours):} $z = \mathrm{clip}(\mu m g / F_{\max},\, 0,\, 1)$, yielding $z \in [0.10, 0.44]$.
    \item \textbf{Mean value:} Fixed $z = 0.245$ (the in-distribution mean of $z$).
    \item \textbf{Fixed at minimum:} $z = 0$.
    \item \textbf{Fixed at maximum:} $z = 1$.
\end{itemize}

\subsection{Hardware Deployment Details}
\label{app:hardware}

\subsubsection{Robot Platform}

All real-world experiments are conducted on a legged mobile manipulator equipped with a 6-DoF arm and a two-finger gripper.
The learned policy is deployed zero-shot without any retraining or fine-tuning.

\subsubsection{Perception Pipeline}

The robot detects the target object using the onboard hand camera.
Object detection uses OWL-ViT~\cite{minderer2022simple} for bounding-box proposal, followed by SAM~\cite{kirillov2023segment} for precise mask extraction.
A pixel on the detected mask edge is passed to the robot's grasp service~\cite{spotsdk2024} for boundary grasp acquisition.
After grasping, the gripper is opened to establish a push-from-edge contact configuration.

\subsubsection{Impedance-Based Force Probing}

The interaction descriptor $z$ is estimated via impedance-based probing.
Table~\ref{tab:probe_params} lists the probing parameters.

\begin{table}[h]
\centering
\caption{Impedance probing parameters.}
\label{tab:probe_params}
\begin{tabular}{ll}
\toprule
\textbf{Parameter}                     & \textbf{Value}         \\
\midrule
Probe step size ($\delta$)             & $0.005$\,m             \\
Maximum probe steps ($N$)              & $30$                   \\
F/T sampling rate                      & $\sim\!50$\,Hz         \\
Sample window per step                 & $1.0$\,s               \\
Force drop threshold ($\Delta f$)      & $4.0$\,N               \\
Consecutive drop steps for breakaway   & $2$                    \\
Probing stiffness                      & $250$\,N/m             \\
Probing damping                        & $30$\,Ns/m             \\
Pre-probe settling time                & $3.0$\,s               \\
$F_{\max}$ (normalizer)               & $200$\,N               \\
\bottomrule
\end{tabular}
\end{table}

\subsubsection{Policy Execution}

Table~\ref{tab:exec_params} lists the parameters governing policy execution on hardware.

\begin{table}[h]
\centering
\caption{Hardware execution parameters.}
\label{tab:exec_params}
\begin{tabular}{ll}
\toprule
\textbf{Parameter}                  & \textbf{Value}         \\
\midrule
Maximum policy steps                & $30$                   \\
Control step duration ($\Delta t_{\mathrm{ctrl}}$) & $2.0$\,s \\
Action scale (admittance gain)      & $0.75$\,m              \\
Yaw correction scale                & $0.05$\,rad            \\
Success progress threshold          & $0.95$                 \\
Deviation tolerance (abort)         & $1.0$\,m               \\
Inter-step settling duration        & $3.0$\,s               \\
\bottomrule
\end{tabular}
\end{table}

\subsubsection{Impedance Controller Gains}

The Cartesian impedance controller gains used during pushing are:
\begin{itemize}
    \item Translational stiffness $K_{\mathrm{eff}}$: $300$\,N/m (planar $x$-$y$), $500$\,N/m (vertical $z$).
    \item Translational damping $\zeta$: $45$\,Ns/m (planar), $45$\,Ns/m (vertical).
    \item Rotational stiffness: $20$\,Nm/rad per axis.
    \item Rotational damping: $1$\,Nms/rad per axis.
\end{itemize}

These gains were tuned experimentally to ensure stable, non-oscillatory contact during pushing without excessive compliance.

\subsubsection{Test Object}

The object used in all hardware experiments is a rigid cardboard box with dimensions $0.465 \times 0.61 \times 0.63$\,m (width $\times$ depth $\times$ height), consistent with the simulation model.
Two payload conditions were tested: $8.8$\,kg and $12.5$\,kg.

\subsubsection{Reference Path}

Hardware experiments use a $60^\circ$ circular arc of radius $1.5$\,m, marked on the floor with red tape.

\subsection{Software and Reproducibility}
\label{app:software}

\subsubsection{Software Versions}

Table~\ref{tab:software} lists the primary software dependencies used for training and evaluation.

\begin{table}[h]
\centering
\caption{Software dependencies.}
\label{tab:software}
\begin{tabular}{ll}
\toprule
\textbf{Package} & \textbf{Version} \\
\midrule
Python     & 3.10+  \\
PyTorch    & 2.7    \\
MuJoCo     & 3.3    \\
Gymnasium  & 1.2    \\
NumPy      & 2.2    \\
SciPy      & 1.15   \\
\bottomrule
\end{tabular}
\end{table}

\subsubsection{Training Seeds and Reproducibility}

Training is conducted across $10$ random seeds (seeds $0$--$9$) with deterministic algorithm settings enabled in PyTorch.
Separate RNG instances with fixed seed offsets control replay buffer sampling, exploration noise, evaluation arc geometry randomization, and target policy noise, ensuring full reproducibility.

For the results reported in the main paper, $5$ seeds are used for computing mean and standard deviation.

\subsubsection{Compute}

Training runs were executed on a single machine. Each seed trains for $3{,}001$ episodes (approximately $150$k--$300$k environment steps depending on episode length).
Training all $10$ seeds to convergence takes approximately 1--3 hours on a consumer-grade laptop GPU (in our case, Apple M2 Pro, using Metal Performance Shaders (MPS)).
